% ODE Inversion Section — to be \input{} into main.tex
% Add to graphicspath: {../results/ode inversion/}

\section{ODE Inversion and Image Editing}

A key advantage of flow matching over stochastic diffusion is that the learned ODE is deterministic and invertible.
Given a real image $x_1$, we recover its latent noise vector $z_0$ by solving the ODE \emph{backwards} from $t{=}1$ to $t{=}0$, then perform controlled editing by manipulating $z_0$ and re-generating.

\textbf{Reconstruction.}
We encode 10 CIFAR-10 test images (one per class) to noise via backward ODE, then decode them back via forward ODE (both using midpoint, NFE=100).
The reconstructions are near-perfect (Figure~\ref{fig:recon}), confirming that the learned velocity field defines an accurate invertible mapping.

\textbf{Latent perturbations.}
Adding Gaussian noise $\epsilon$ to the recovered $z_0$ at increasing scales $\sigma$ produces a smooth degradation: $\sigma{=}0.1$ yields barely perceptible changes, while $\sigma{=}1.5$ destroys the original content entirely (Figure~\ref{fig:perturb}).
This confirms the latent space is smooth and locally stable.

\textbf{Real image interpolation.}
We encode pairs of real images to their latent vectors $z_0^A, z_0^B$ and decode linear blends $(1{-}\alpha)\,z_0^A + \alpha\,z_0^B$.
The resulting paths show semantically meaningful transitions between real images (Figure~\ref{fig:real_interp}), unlike random noise interpolation which cannot guarantee passing through real images at the endpoints.

\textbf{Attribute editing.}
We compute class-mean latent vectors from 50 images per class and use direction vectors $\bar{z}_0^{\text{target}} - \bar{z}_0^{\text{source}}$ to edit images.
Applying the direction with varying strength $\lambda$ produces gradual semantic shifts (e.g., horse$\to$car, cat$\to$dog) as shown in Figure~\ref{fig:attr_edit}.
At 32$\times$32 resolution, the effects are subtle but clearly directional.

\begin{figure}[t]
\centering
\includegraphics[width=\columnwidth]{{../results/ode inversion/reconstruction}.png}
\caption{ODE inversion reconstruction. Top: originals. Middle: reconstructed via encode$\to$decode. Bottom: pixel difference ($5\times$ amplified). Near-black differences confirm accurate inversion.}
\label{fig:recon}
\end{figure}

\begin{figure}[t]
\centering
\includegraphics[width=\columnwidth]{{../results/ode inversion/perturbations}.png}
\caption{Latent perturbation: $z_0 + \sigma\epsilon$ decoded at increasing $\sigma$. Small perturbations preserve identity; large ones destroy it.}
\label{fig:perturb}
\end{figure}

\begin{figure}[t]
\centering
\includegraphics[width=\columnwidth]{{../results/ode inversion/attribute_editing}.png}
\caption{Attribute editing via class direction vectors at varying $\lambda$. The source image ($\lambda{=}0$) is gradually pushed toward the target class semantics.}
\label{fig:attr_edit}
\end{figure}

\begin{figure}[t]
\centering
\includegraphics[width=\columnwidth]{{../results/ode inversion/inversion_trajectory}.png}
\caption{Inversion trajectory: real images dissolve into Gaussian noise as the ODE is solved backwards from $t{=}1$ to $t{=}0$.}
\label{fig:real_interp}
\end{figure}
